% !PRINTED 564
\weberpage{564}
%<<CONT
also sind $\xi$, $\xi_2$, $\xi_4$ die Wurzeln der cubischen Gleichung
\weq{16}{x^3 - \zeta x^2 - x - 1 + \zeta = 0.}

\begin{center}
\S.\,169.

\gesperrt{Zur\"uckf\"uhrung der Kreistheilungsgleichung auf}\\
\gesperrt{reine Gleichungen.}
\end{center}

\gesperrt{Gauss} hat schon in seiner ersten Darstellung in den disq.
arithm. die Kreistheilungsgleichungen durch Benutzung der Resolventen
direct auf \gesperrt{reine Gleichungen} zur\"uckgef\"uhrt\footnote{\gesperrt{Gauss}, disq. arithm. art. 359, 360; disq. circa aequationes puras ulterior evolutio, Werke Bd. II. \gesperrt{Lagrange}, r\'es. des \'equations num\'eriques. \gesperrt{Jacobi}, ,,Ueber die Kreistheilung und ihre Anwendung in der Zahlentheorie``. Werke Bd. 6. \gesperrt{Kummer}, Crelle's Journal, Bd. 35 und Abhandl. d. Berl. Akademie 1856. Zu erw\"ahnen sind noch \gesperrt{Eisenstein} und \gesperrt{Cauchy}. Die Lehre von der Kreistheilung ist im Zusammenhange dargestellt und von den historischen Nachweisen begleitet in dem Buche von \gesperrt{Bachmann}, ,,Die Lehre von der Kreistheilung``. Leipzig 1872.}. Um
dies durchzuf\"uhren, setzen wir
\[ n - 1 = m, \]
und bezeichnen, wie im vorigen Abschnitt, mit $\varepsilon$ eine primitive
Einheitswurzel $m^{\mathrm{ten}}$ Grades, mit $r, r_1, r_2 \ldots r_{m-1}$ die $n^{\mathrm{ten}}$ Einheitswurzeln
 in der oben festgesetzten Reihenfolge. Die \gesperrt{Lagrange'}schen
Resolventen sind dann
\[ (\varepsilon^\lambda, r) = \sum_{0,\,n-2}^{h} \varepsilon^{\lambda h} r_h = \sum_{0,\,n-2}^{h} \varepsilon^{\lambda h} r^{g^h}, \]
worin $\lambda$ ein beliebiger, nach dem Modul $m$ genommener Exponent
 ist.
Setzen wir $g^h = s$, also $h \equiv \mathrm{ind}\,s \pmod{m}$; so k\"onnen wir
auch setzen
\weq{1}{(\varepsilon^\lambda, r) = \sum_{1,\,n-1}^{s} \varepsilon^{\lambda\,\mathrm{ind}\,s} r^s.}

Als Grundlage f\"ur die weitere Reduction dieser Ausdr\"ucke
dient die Berechnung des Productes zweier solcher Resolventen.

Wir bezeichnen mit $\mu$ eine zweite Zahl wie $\lambda$ und verstehen
unter $s$ und $t$ zwei Zeichen, die von einander unabh\"angig je
ein Restsystem nach dem Modul $n$ durchlaufen, mit Ausschluss
der Null.
Dann ist
\[ (\varepsilon^\lambda, r)(\varepsilon^\mu, r) = \sum^{s} \sum^{t} r^{(s+t)} \varepsilon^{\lambda\,\mathrm{ind}\,s + \mu\,\mathrm{ind}\,t}. \]
